% Copyright 2026 icekylinx and GPT-6 Astra Pro. Apache-2.0.
% Curated from sat_research_checkpoint.zip: paper/algebraic_streaming_coefficients.tex
% Complete finite proof retained; no SAT dependency or formal receipts imported.
\subsection{Computing the integral crossing coefficients without large tables}
\label{sec:algebraic-streaming-coefficients}

The following finite algorithm supplies the construction cost needed
by the algebraic circuit consequences.  It uses only the modes and
the integral Schur family already defined above; it has no
tensor-network or satisfiability-algorithm premise.

\begin{lemma}[Explicit computation of integral crossing coefficients]
\label{lem:algebraic-coefficient-computation}
For the weight parameters $h,E\ge1$ of
Theorem~\ref{thm:integral-crossing}, any entry of a crossing piece
or of either of its displayed factorizations can be computed in
$2^{O(\sqrt{h+E})}$ bit operations and
$\operatorname{poly}(h+E)$ space.  All rational intermediates and
the final integer have bit length polynomial in $h+E$.
The implied constants are absolute.
\end{lemma}

\begin{proof}
Every state, input, output, or intermediate partition in the
crossing has weight $O(h+E)$: in particular, the intermediate
weights $E+s-h+i$ and $E+h-s-i$ are at most $2E+2h$ when
$|s|\le E+h$ and $0\le i\le2h$.  Put $W=4(h+E+1)$ to bound
all these weights.  Partitions of weight at most $W$ can be
enumerated in polynomial space, and their number is
$2^{O(\sqrt W)}$.  The latter bound follows directly from Euler's
product evaluated at $\exp(-1/\sqrt W)$.

\paragraph{A power-sum mode entry.}
Write $m_\sigma(j)$ for the multiplicity of part $j$ in a
partition $\sigma$.  Consider first
\[
 [p_\tau]a[r]p_\sigma.
\]
Return zero unless $r=|\tau|-|\sigma|$.  Normal ordering means
that differentiations act before creations.  Enumerate derivative
multiplicities $0\le d_j\le m_\sigma(j)$.  The required creation
multiplicities are forced to be
\[
 c_j=m_\tau(j)-m_\sigma(j)+d_j.
\]
Discard a vector with a negative $c_j$; otherwise
$0\le c_j\le m_\tau(j)$.  The total mode-shift condition is then
automatic.  The differentiation contributes the integer
\[
 D(d)=\prod_{j\ge1}j^{d_j}(m_\sigma(j))_{d_j},
\]
where $(m)_d=m(m-1)\cdots(m-d+1)$.

For a monomial state $a=p_\mu=p_{\mu_1}\cdots p_{\mu_\ell}$,
form the ordinary homogeneous polynomial
\[
 H_\mu(Y)=\prod_{t=1}^{\ell}
       \left(\sum_{k\ne0}\binom{-k-1}{\mu_t-1}Y_k\right).
\]
In this expression retain only positive $k$ in
$\operatorname{supp}(\sigma)$ and negative $k$ with
$-k\in\operatorname{supp}(\tau)$.  No other differentiations can
survive, and no other creations can occur in the requested output.
The contribution of $(d,c)$ is $D(d)$ times the coefficient of
\[
 \prod_jY_j^{d_j}Y_{-j}^{c_j}
\]
in $H_\mu$.  It is zero unless $\sum_j(d_j+c_j)=\ell$.

For an ordinary homogeneous polynomial $H$ of degree $|e|$,
rectangular finite differences extract precisely this coefficient:
\begin{equation}
\label{eq:algebraic-rectangular-differences}
 [Y^e]H
 =\frac1{\prod_k e_k!}
   \sum_{0\le z_k\le e_k}
       (-1)^{|e|-|z|}\prod_k\binom{e_k}{z_k}H(z).
\end{equation}
Indeed every other monomial of degree $|e|$ has an exponent below
its target in at least one coordinate, and is annihilated by that
coordinate's finite difference.  The same identity works with
coefficients in a univariate polynomial ring.

A derivative vector is a subpartition of $\sigma$, so the number
of such vectors is at most
\[
 \prod_j(m_\sigma(j)+1)\le\sum_{v=0}^W p(v)
 =2^{O(\sqrt W)}.
\]
For each vector, the finite-difference grid has at most
\[
 \prod_j(d_j+1)(c_j+1)
 \le\prod_j(m_\sigma(j)+1)(m_\tau(j)+1)
 =2^{O(\sqrt W)}
\]
points.  A fixed product of these bounds remains
$2^{O(\sqrt W)}$.  At a grid point the displayed product of
$\ell\le W$ short linear forms is evaluated directly.

\paragraph{The transformed state and its weight marker.}
For $a=T(p_\mu)$, use instead
\[
 H_\mu(Y,Z)=
 \prod_{t=1}^{\ell}(-1)^{\mu_t-1}
 \left(\sum_{j=1}^{\mu_t}\binom{\mu_t}{j}Z^j
       \sum_{k\ne0}\binom{-k-1}{j-1}Y_k\right).
\]
The global product of signs is retained.  This polynomial is still
homogeneous of degree $\ell$ in $Y$, and the marker $Z$ records
the weight of each homogeneous state component.  Apply
\eqref{eq:algebraic-rectangular-differences} while storing a
univariate array of degree at most $h$.  For
$a^{\langle i\rangle}$, apply after coefficient extraction the
linear functional
\[
 Z^w\longmapsto\binom{h+w}{i}.
\]
For $a$ itself, sum the marker coefficients.  This avoids expanding
$T(p_\mu)$ into all of its monomials.  It also handles the
inhomogeneous transform without confusing state weight with the
mode shift $r$.

All generalized binomial coefficients have arguments of magnitude
$O(W)$ and bit length $O(W)$.  Products contain at most $W$
factors, and the factorials and falling factorials above have
polynomial bit length.  The univariate arrays have polynomial
rank.  Thus every evaluation, accumulation, and exact division in
this part uses polynomial space and polynomial-bit integers.

\paragraph{Streaming the Schur basis changes.}
For partitions $\lambda,\mu$ of weight $w\le W$, put
$e_j=m_\mu(j)$ and set power sums outside
$\operatorname{supp}(\mu)$ to zero.  Then
\[
 [p_\mu]s_\lambda
 =\frac1{\prod_j e_j!}
   \sum_{0\le z_j\le e_j}
   (-1)^{\sum_j(e_j-z_j)}
   \prod_j\binom{e_j}{z_j}\,s_\lambda(p_j=z_j).
\]
Here weighted homogeneity suffices: a different monomial of
weight $w$ cannot be componentwise at least the target, because
any strict increase would have weight greater than $w$.
Therefore it has a below-target coordinate and vanishes under
that finite difference.  The grid has
$\prod_j(e_j+1)=2^{O(\sqrt W)}$ points.

At each grid point compute the complete functions by
\[
 m h_m=\sum_{j=1}^m p_jh_{m-j},\qquad h_0=1,
\]
and evaluate the Jacobi--Trudi determinant for $s_\lambda$ using
polynomial-time rational arithmetic.  Factorial common
denominators bound the bit lengths polynomially.  Hall
orthogonality gives the reverse coefficient
\[
 [s_\lambda]p_\mu=z_\mu[p_\mu]s_\lambda,
 \qquad z_\mu=\prod_jj^{m_\mu(j)}m_\mu(j)!.
\]

To compute a Schur mode entry, stream over the power-sum
partitions for its state, input, and output, multiply the three
basis-conversion coefficients by the power-sum mode entry just
computed, and accumulate.  For a transformed Schur state use
linearity, $T(s_\lambda)=\sum_\mu[p_\mu]s_\lambda\,T(p_\mu)$,
and the marker procedure.  Each of the three partition sets has
size $2^{O(\sqrt W)}$, so this fixed number of nested enumerations
and coefficient computations still takes $2^{O(\sqrt W)}$ time.
The rational basis coefficients have common denominators dividing
weight factorials; during their finite-difference computation a
fixed power of $W!$ is a sufficient common denominator.  Thus
accumulating exponentially many terms does not create
exponential bit lengths.

Finally, obtain a factorization entry from its indicated mode
entry and sign.  To obtain an entry of a full crossing piece,
stream its polynomially many $i$ values and the intermediate
partitions.  This is only a further fixed number of the same
enumerations.  Theorem~\ref{thm:integral-crossing} proves that
the final coordinates are integers.  Any final exact division
is performed over the integers before specialization to a target
ring.  No large table is materialized: partition indices, the
derivative vector, one finite-difference point, polynomial-size
univariate arrays, and accumulators suffice.  This proves all
three stated bounds.
\end{proof}
